\documentclass[12pt,a4paper,oneside]{report}

% ---------------------------------------------------------
% PACCHETTI
% ---------------------------------------------------------

\usepackage[utf8]{inputenc}
\usepackage[T1]{fontenc}
\usepackage[italian]{babel}

\usepackage{amsmath}
\usepackage{amssymb}

\usepackage{graphicx}
\usepackage{float}

\usepackage{booktabs}
\usepackage{array}

\usepackage{listings}
\usepackage{xcolor}

\usepackage{geometry}
\usepackage{setspace}

\usepackage{hyperref}
\usepackage{url}

\usepackage{caption}
\usepackage{subcaption}

% ---------------------------------------------------------
% IMPOSTAZIONE PAGINA
% ---------------------------------------------------------

\geometry{
    a4paper,
    left=3cm,
    right=3cm,
    top=3cm,
    bottom=3cm
}

\onehalfspacing

% ---------------------------------------------------------
% HYPERREF
% ---------------------------------------------------------

\hypersetup{
    colorlinks=true,
    linkcolor=black,
    citecolor=black,
    urlcolor=blue,
    pdftitle={Titolo della tesi},
    pdfauthor={Nome Cognome}
}

% ---------------------------------------------------------
% CODICE
% ---------------------------------------------------------

\lstset{
    basicstyle=\ttfamily\small,
    breaklines=true,
    frame=single,
    numbers=left,
    numberstyle=\tiny,
    showstringspaces=false
}

% ---------------------------------------------------------
% INFORMAZIONI TESI
% ---------------------------------------------------------

\newcommand{\thesistitle}{gestione della compra/vendita di energia }
\newcommand{\authorname}{Daniele Quattrociocchi}
\newcommand{\degree}{Laurea in Informatica}
\newcommand{\supervisor}{Prof. Enrico Tronci}
\newcommand{\academicyear}{2025--2026}

% ---------------------------------------------------------
% DOCUMENTO
% ---------------------------------------------------------

\begin{document}

% =========================================================
% FRONTESPIZIO
% =========================================================

\begin{titlepage}

    \centering

    % Inserire qui eventualmente il logo Sapienza
    % \includegraphics[width=0.25\textwidth]{images/sapienza-logo.png}

    \vspace*{1cm}

    {\Large \textbf{SAPIENZA UNIVERSITÀ DI ROMA}\par}

    \vspace{1cm}

    {\large Facoltà di Ingegneria dell'Informazione, Informatica e Statistica\par}

    \vspace{0.5cm}

    {\large Corso di Laurea in Informatica\par}

    \vspace{2.5cm}

    {\LARGE \textbf{\thesistitle}\par}

    \vspace{1.5cm}

    {\large Tesi di Laurea\par}

    \vfill

    \begin{flushleft}
        \large
        \textbf{Relatore:} \supervisor

        \vspace{0.5cm}

        \textbf{Candidato:} \authorname
    \end{flushleft}

    \vspace{1cm}

    {\large Anno Accademico \academicyear\par}

\end{titlepage}


% =========================================================
% FRONT MATTER
% =========================================================

\pagenumbering{roman}

% ---------------------------------------------------------
% ABSTRACT
% ---------------------------------------------------------

\chapter*{Abstract}

% L'abstract verrà scritto alla fine.
% Massimo una pagina.

Inserire qui l'abstract della tesi.

\clearpage


% ---------------------------------------------------------
% INDICE
% ---------------------------------------------------------

\tableofcontents

\clearpage


% =========================================================
% MAIN MATTER
% =========================================================

\pagenumbering{arabic}


% =========================================================
% 1. INTRODUCTION
% =========================================================

\chapter{Introduzione}

% Nota:
% Il professore suggerisce di scrivere questo capitolo alla fine.

\section{Contesto}

Descrizione del contesto generale del problema.

\section{Motivazioni}

Motivazioni che hanno portato allo sviluppo del progetto.

\section{Contributi}

Descrizione dei contributi apportati dal lavoro.

\section{Stato dell'arte}

Descrizione dello stato dell'arte.

% Esempio di citazione:
% \cite{paper_example}

\section{Struttura della tesi}

La tesi è organizzata come segue.
Il Capitolo~\ref{chap:background} introduce il background necessario
alla comprensione del lavoro.
Il Capitolo~\ref{chap:methods} descrive i metodi utilizzati.
Il Capitolo~\ref{chap:implementation} descrive l'implementazione.
Il Capitolo~\ref{chap:results} presenta i risultati sperimentali.
Infine, il Capitolo~\ref{chap:conclusions} presenta le conclusioni
e i possibili sviluppi futuri.


% =========================================================
% 2. BACKGROUND
% =========================================================

\chapter{Contesto teorico}
\label{chap:background}

% Inserire qui tutto il background necessario per comprendere
% il progetto.

\section{Argomento teorico 1}

Descrizione del primo concetto necessario.

\section{Argomento teorico 2}

Descrizione del secondo concetto necessario.

\section{Argomento teorico 3}

Descrizione di eventuali algoritmi, tecnologie o strumenti.


% =========================================================
% 3. METHODS
% =========================================================

\chapter{Metodologia}
\label{chap:methods}

% Descrizione della progettazione del sistema.
% NON inserire dettagli di implementazione o codice.

\section{Definizione del problema}

Definizione formale o informale del problema affrontato.

\section{Approccio proposto}

Descrizione dell'approccio proposto.

\section{Progettazione del sistema}

Descrizione della progettazione generale del sistema.

\section{Algoritmi}

Descrizione degli algoritmi utilizzati o progettati.

% Esempio:
%
% \begin{equation}
%     ...
% \end{equation}


% =========================================================
% 4. IMPLEMENTATION
% =========================================================

\chapter{Implementazione}
\label{chap:implementation}

% Descrizione dell'implementazione concreta.
% Qui possono essere utilizzati pseudocodice e diagrammi.

\section{Architettura del sistema}

Descrizione dell'architettura del sistema.

% Esempio di figura:
%
% \begin{figure}[H]
%     \centering
%     \includegraphics[width=0.8\textwidth]{images/architecture.png}
%     \caption{Architettura del sistema.}
%     \label{fig:architecture}
% \end{figure}

\section{Componenti principali}

Descrizione dei principali componenti software.

\section{Dettagli implementativi}

Descrizione delle principali scelte implementative.

\section{Implementazione degli algoritmi}

Descrizione dell'implementazione degli algoritmi.

% Esempio di codice:
%
% \begin{lstlisting}[language=C++, caption={Esempio di codice}]
% int main() {
%     return 0;
% }
% \end{lstlisting}


% =========================================================
% 5. EXPERIMENTAL RESULTS
% =========================================================

\chapter{Risultati sperimentali}
\label{chap:results}

% Questa dovrebbe essere una delle sezioni più importanti
% della tesi.

\section{Obiettivi}

L'attività sperimentale ha l'obiettivo di valutare l'efficacia delle strategie
di compravendita implementate nel simulatore. In particolare, il confronto
principale riguarda la strategia \texttt{smartGeometricChoise} e la strategia
\texttt{randomChoice}, utilizzata come baseline.

Gli esperimenti sono stati progettati per:

\begin{itemize}
    \item misurare il rendimento medio prodotto dalle due strategie;
    \item analizzare la variabilità dei risultati tramite la deviazione standard;
    \item studiare l'influenza dei principali parametri della strategia Smart;
    \item verificare in quali configurazioni la strategia Smart migliora il
          comportamento della baseline casuale.
\end{itemize}

Il rendimento medio non viene quindi considerato come unico criterio di
valutazione: una strategia più redditizia, ma anche più variabile, deve essere
interpretata tenendo conto del corrispondente aumento dell'incertezza.


% ---------------------------------------------------------
% 5.2 SETTING
% ---------------------------------------------------------

\section{Configurazione sperimentale}

Descrizione dell'ambiente sperimentale.

\subsection{Hardware}

Descrizione dell'hardware utilizzato.

\subsection{Software}

Descrizione del software, delle librerie e delle versioni utilizzate.

\subsection{Parametri}

I test sono stati eseguiti mantenendo fissi la durata della simulazione, il
budget per operazione, il saldo iniziale e la soglia di vendita. Sono stati
variati la finestra geometrica, la percentuale accantonata e la soglia di
acquisto, secondo i valori riportati in Tabella~\ref{tab:experimental-parameters}.

\begin{table}[H]
    \centering
    \caption{Parametri utilizzati nella campagna sperimentale.}
    \label{tab:experimental-parameters}
    \begin{tabular}{ll}
        \toprule
        Parametro & Valori considerati \\
        \midrule
        Finestra simulazione & 300 \\
        Budget per operazione & 10 \\
        Saldo iniziale & 1000 \\
        Finestra geometrica & 10, 20, 50, 100 \\
        Set Aside (\%) & 10, 20, 30, 50 \\
        Buy Threshold & 0.9, 0.8, 0.7 \\
        Sell Threshold & 1.0 \\
        \bottomrule
    \end{tabular}
\end{table}

La strategia \texttt{randomChoice} non utilizza la finestra geometrica né le
soglie di acquisto e vendita per prendere le decisioni: per ogni rete e per
ogni istante seleziona casualmente l'azione da eseguire. Per questo motivo non
è metodologicamente corretto associare una singola riga Smart a una riga
Random avente gli stessi valori di tutti i parametri. Random viene invece
trattata come baseline, mentre i parametri specifici della strategia Smart
sono analizzati separatamente.


% ---------------------------------------------------------
% 5.3 CASE STUDIES
% ---------------------------------------------------------

\section{Casi di studio}

I casi di studio sono costruiti a partire dalle combinazioni dei parametri
riportati nella sezione precedente. Per ogni configurazione vengono raccolti
il rendimento medio (\textit{Average Yield}), la deviazione standard e il
guadagno finale medio.

\subsection{Baseline casuale}

La strategia \texttt{randomChoice} costituisce il riferimento sperimentale.
Nel campione analizzato produce un rendimento medio compreso tra il 7,40\%
e il 7,58\%, una deviazione standard prossima al 12\% e un guadagno finale
medio compreso tra circa 74 e 76 unità monetarie. La variazione residua tra
le configurazioni è compatibile con il carattere casuale della strategia e
con l'effetto della percentuale accantonata.

\subsection{Strategia Smart}

Per \texttt{smartGeometricChoise} sono state analizzate quattro finestre
geometriche e tre soglie di acquisto. La percentuale accantonata è stata
considerata come parametro secondario, poiché in diverse configurazioni il
suo effetto sul rendimento risulta nullo o molto contenuto.


% ---------------------------------------------------------
% 5.4 CORRECTNESS
% ---------------------------------------------------------

\section{Correttezza}

Descrizione degli esperimenti utilizzati per verificare la correttezza
dell'implementazione.

\subsection{Test di correttezza}

Descrizione dei test effettuati.

\subsection{Risultati}

Presentazione dei risultati relativi alla correttezza.


% ---------------------------------------------------------
% 5.5 COMPUTATIONAL EVALUATION
% ---------------------------------------------------------

\section{Valutazione computazionale}

Valutazione delle prestazioni computazionali.

\subsection{Tempo di esecuzione}

Analisi dei tempi di esecuzione.

\subsection{Utilizzo della memoria}

Analisi dell'utilizzo della memoria.

\subsection{Utilizzo della CPU}

Analisi dell'utilizzo della CPU.


% ---------------------------------------------------------
% 5.6 TECHNICAL EVALUATION
% ---------------------------------------------------------

\section{Valutazione tecnica}

La valutazione tecnica considera il rendimento e la sua variabilità per
confrontare la strategia Smart con il comportamento della baseline casuale.

\subsection{Risultati sperimentali}

I risultati mostrano una forte dipendenza della strategia Smart dalla soglia
di acquisto. Considerando, a titolo rappresentativo, una finestra geometrica
pari a 20 e una percentuale accantonata pari al 10\%, si osserva il seguente
andamento:

\begin{table}[H]
    \centering
    \caption{Risultati rappresentativi per finestra geometrica pari a 20.}
    \label{tab:smart-representative-results}
    \resizebox{\textwidth}{!}{%
    \begin{tabular}{lrrr}
        \toprule
        Buy Threshold & Average Yield (\%) & Deviazione standard (\%) & Guadagno medio \\
        \midrule
        0.9 & 5.64 & 9.83 & 56.36 \\
        0.8 & 8.23 & 11.76 & 82.33 \\
        0.7 & 10.77 & 13.21 & 107.69 \\
        \bottomrule
    \end{tabular}%
    }
\end{table}

Rispetto alla baseline, il rendimento della strategia Smart passa da valori
inferiori a Random con soglia 0.9 a valori superiori con soglia 0.8 e 0.7.
La configurazione più redditizia osservata nel campione è quella con finestra
geometrica pari a 10 e Buy Threshold pari a 0.7: raggiunge un rendimento
medio del 12,25\%, contro circa il 7,4\% della baseline.

\subsection{Confronto}

Il confronto viene effettuato mantenendo costanti i parametri comuni e
considerando Random come valore di riferimento. Poiché Random non dipende
dalla finestra geometrica e dalle soglie usate da Smart, il risultato corretto
non è un confronto uno-a-uno tra configurazioni omonime, ma un confronto tra
la baseline e l'insieme delle configurazioni Smart.

La Tabella~\ref{tab:comparison-summary} riassume tre configurazioni Smart e
la fascia di risultati osservata per Random.

\begin{table}[H]
    \centering
    \caption{Confronto sintetico tra Random e alcune configurazioni Smart.}
    \label{tab:comparison-summary}
    \resizebox{\textwidth}{!}{%
    \begin{tabular}{lrrr}
        \toprule
        Configurazione & Average Yield (\%) & Deviazione standard (\%) & Guadagno medio \\
        \midrule
        Random (intervallo osservato) & 7.40--7.58 & 11.95--11.99 & 74.02--75.75 \\
        Smart, geom. 20, buy 0.9 & 5.64 & 9.83 & 56.36 \\
        Smart, geom. 20, buy 0.8 & 8.23 & 11.76 & 82.33 \\
        Smart, geom. 20, buy 0.7 & 10.77 & 13.21 & 107.69 \\
        \bottomrule
    \end{tabular}%
    }
\end{table}

\subsection{Discussione}

Il confronto evidenzia che la strategia Smart non è uniformemente superiore
alla strategia Random. Con Buy Threshold pari a 0.9, la baseline casuale
ottiene un rendimento maggiore; riducendo la soglia a 0.8, Smart supera
invece il rendimento medio di Random; con soglia 0.7 il vantaggio diventa più
marcato.

Anche la finestra geometrica influenza il risultato. Con Buy Threshold pari
a 0.7, il rendimento passa dal 12,25\% con finestra 10 al 10,77\% con
finestra 20, all'8,91\% con finestra 50 e al 7,37\% con finestra 100. Nel
campione considerato, finestre più brevi sono quindi associate a rendimenti
maggiori. Questo risultato è compatibile con una maggiore reattività alle
variazioni recenti dei prezzi, ma tale interpretazione deve essere riferita
alla configurazione e ai dati analizzati, senza essere generalizzata oltre
l'esperimento.

Il rendimento maggiore è accompagnato da una maggiore variabilità: la
configurazione Smart migliore raggiunge una deviazione standard del 14,63\%,
contro circa il 11,95\% osservato per Random. La scelta della strategia non
dovrebbe quindi basarsi esclusivamente sul rendimento medio, ma sul rapporto
tra rendimento ottenuto e variabilità dei risultati.


% =========================================================
% 6. CONCLUSIONS
% =========================================================

\chapter{Conclusioni}
\label{chap:conclusions}

Riassunto del lavoro svolto.

\section{Sintesi}

Sintesi degli obiettivi e del lavoro realizzato.

\section{Sviluppi futuri}

Possibili sviluppi futuri del progetto.


% =========================================================
% REFERENCES
% =========================================================

\bibliographystyle{plain}
\bibliography{references}

\end{document}